\documentclass{article}
\usepackage[margin=1in]{geometry}
\usepackage{mathtools, amsfonts, amsthm, xcolor, listings}

\title{HW 03}
\author{Sam Ly}

\begin{document}
\maketitle

\section*{Exercise 1}
\noindent
Tools and resources used: 

\begin{enumerate}
  \item[Exercise 1]
    {\color{red} Points attempted: 2}.
    {\color{blue} Time spent: 5 minutes}.
    {\color{brown} Difficulty: 1/10}.
    {\color{green!50!black} Self-assessed score: 1}.
    {\color{magenta} Questions/Feedback: None.}
  \item[Choice Exercise 2]
    {\color{red} Points attempted: 2}.
    {\color{blue} Time spent: 5 minutes}.
    {\color{brown} Difficulty: 1/10}.
    {\color{green!50!black} Self-assessed score: 2}.
    {\color{magenta} Questions/Feedback: None.}
  \item[Choice Exercise 3]
    {\color{red} Points attempted: 6}.
    {\color{blue} Time spent: 30 minutes}.
    {\color{brown} Difficulty: 3/10}.
    {\color{green!50!black} Self-assessed score: 6}.
    {\color{magenta} Questions/Feedback: Mathematical correctness, clarity.}
\end{enumerate}
\pagebreak 

\section*{Exercise 2}
Done! I discussed my solution to HW02 Exercise 9.4.
\pagebreak

\section*{Exercise 3}
\begin{enumerate}
    \item {`'
        Let \((G, *) = (\mathbb{Z}, +)\), and let \(a = 1 \in G\) and \(b = 3 \in G\). Compute \(a^4 * b^{-2} \in \mathbb{Z}\).
      
        We see that \(a^4 * b^{-2}\) expands to \(1+1+1+1-3-3 = -2\).
    }
    \item {
        Let \(g \in G\), and let \(\langle g \rangle\) be the cyclic subgroup of \(G\) generated by \(g\). Prove that the order of the element \(g\) is equal to the order of the subgroup \(\langle g \rangle\).
    
        \begin{proof}
            First, we notice that the order of \(\langle g \rangle\) is countably
            infinite. Thus, traditional counting techniques may not be used.

            We can rewrite all elements of \(\langle g \rangle\) as \(\{g^n \mid n \in \mathbb{Z}\}\).
            We can also rewrite the subgroup generated by \(g\) as \(\{g^n \mid n \in \mathbb{Z}\}\).

            Therefore, these two group are exactly the same, and have exactly the same order.
        \end{proof}
    }
    \item {
        Consider the cyclic group \(C_{20}\), where \(r\) is an \(18^\circ\) rotation, and thus \(|r| = 20\). Find all elements \(r^n \in C_{20}\) such that \(\langle r^n \rangle = C_{20}\).
    
        For all \(n \in \{1, 3, 7, 9, 11, 13, 17, 19\}\), \(\langle r^n \rangle = C_{20}\).
    }
    \item {
        Give an example of an infinite group \((G,*)\) and a subset \(H \subseteq G\) which is closed under \(*\), but where \(H\) is not a subgroup of \(G\).

        Let \(H = 2\mathbb{Z} + 1\). \(H\) is closed under \(*\) but does not include 
        an identity element, so it is not a subgroup of \(G\).
    }
    \item {
        Prove that if \(a \equiv a' \pmod n\), \(b \equiv b' \pmod n\), and \(a + b \equiv c \pmod n\), then \[         a' + b' \equiv c \pmod n.       \]

        \begin{proof}
            By definition, \(a - a' = kn\) and \(b = b' = ln\). If \(a + b \equiv c \pmod n\),
            then \(a + b - c = jn\). 

            With some algebra, we find 
            \[a = kn + a', b = ln + b'\].

            Then, we write \(kn + a' + ln + b' - c = jn\). We can then subtract over 
            some terms, then factor to get
            \[a' + b' - c = jn - kn -ln\]
            \[a' + b' - c = (j - k - l)n\].
            
            Therefore, \(a' + b' \equiv c \pmod n\). 
        \end{proof}
    }
\end{enumerate}
\pagebreak

\section*{Choice Exercise 6 [15]}
\begin{enumerate}
    \item {
    [4]
        Prove that the set of \(3 \times 3\) matrices that have one \(\pm 1\) per row and column forms a group under matrix multiplication. This group is called the octahedral group, sometimes written \(B_3\). Without looking it up, determine the order of this group.

        \begin{proof}
            To prove that a given set of matrices forms a group under matrix multiplication,
            we must prove that it satisfies the group axioms. 

            First, associativity is inherited from matrix multiplication.
            Then, we see that the set of \(3 \times 3\) matrices that have one 
            \(\pm 1\) per row and column includes the identity matrix \(I\). 

            Also, we see that this set of matrices must have linearly independent 
            columns, because each of the columns has an element in a unique row. 
            Thus, the matrices are invertible. 

            The most difficult part of this proof is proving that this is set of 
            matrices is closed under matrix multiplication. First, let \(a, b \in B_3\).
            Let the 3 rows of \(a\) be the vectors \(a^1, a^2, a^3\) (this 
            non-standard notation, however it is helpful for this proof), and the 3
            columns of \(b\) be the vectors \(b_1, b_2, b_3\). 

            Notice that having ``one \(\pm 1\) per row and column'' implies that 
            the the vectors \(a^1, a^2, a^3\) each must have its non-zero element 
            be in a unique index. That is, if \(a^1 = (1,0,0)\), then \(a^2, a^3\) can't 
            be \((1,0,0)\) or \((-1,0,0)\). This same property applies to \(b_1, b_2, b_3\).

            Now, let \(c = a \times b\). I will denote each element of \(c\) as 
            \(c^i_j\), for \(i, j \in \{1, 2, 3\}\), where \(i\) is the row index 
            and \(j\) is the column index. By definition of matrix multiplication,
            \(c^i_j = a^i \cdot b_j\) (dot product). The only time this value is 
            non-zero is if the the index of the non-zero element in \(a^i\) and \(b_j\) match. 

            Because the index of the non-zero element of the rows of \(a\) and 
            the columns of \(b\) are unique, if \(c^i_j = a^i \cdot b_j \not= 0\), 
            all other elements of \(c\) in the same row/col of \(c^i_j\) must be 0.
            
            Therefore, the set of \(3 \times 3\) matrices that have one \(\pm 1\) 
            per row and column is closed under matrix multiplication, so, they form 
            a group \(B_3\).
        \end{proof}
    }

    \item {
        Done!

        \begin{lstlisting}
module blob() {
  cylinder(20,10,5);
  translate([-5,0,10]) rotate([45,0,0]) cube([20,10,5]);
}

blob();
        \end{lstlisting}
    }
\end{enumerate}
\pagebreak

\end{document}
